%------------------------------------------------------------------------------%
\documentclass[fleqn,utf8,aspectratio=169,12pt]{beamer}

\usepackage{gaal}

\title{Módulo, Soma e Produto por Escalar}
\subtitle{Operações com  Vetores}

%------------------------------------------------------------------------------%
\begin{document}

\addtitle

%------------------------------------------------------------------------------%
\section{Módulo de um Vetor}

%------------------------------------------------------------------------------%
\begin{frame}{Módulo de um Vetor}
  \emph{Módulo}, \emph{Comprimento} ou \emph{Norma} do vetor

  \pause
  é o comprimento do segmento orientado
  \pause
  \[
    \norm{\vec{u}}
    \pause
    = \abs{\vec{u}}
    \pause
    = \abs{\overrightarrow{AB}}
  \]

  \pause
  Direção e sentido não interferem no módulo

  \pause
  O módulo é uma grandeza escalar (número real) não negativa
  \[
    \pause
    \norm{\vec{u}} \in \R
    \qquad\qquad
    \norm{\vec{u}} \geq 0
  \]
\end{frame}

%------------------------------------------------------------------------------%
\framedgraphic*{Interpretação Geométrica}{E-02-Operacoes-modulo-1}{}
\framedgraphic*{Interpretação Geométrica}{E-02-Operacoes-modulo-2}{}
\framedgraphic*{Interpretação Geométrica}{E-02-Operacoes-modulo-3}{}

%------------------------------------------------------------------------------%
\begin{frame}[t]{Cálculo do Módulo}
  \vspace{1\baselineskip}

  Para um vetor no plano $\vec{u} \in \R^2$
  \[
    \pause
    \vec{u}
    \pause
    = ( u_1, u_2 )
    \pause
    = \vetor{u_1}{u_2}
  \]

  \pause
  Pelo Teorema de Pitágoras
  \[
    \pause
    \norm{\vec{u}}
    \pause
    = \sqrt{ u_1^2 + u_2^2 }
  \]
\end{frame}

%------------------------------------------------------------------------------%
\begin{frame}[t]{Cálculo do Módulo}
  \vspace{1\baselineskip}

  Para um vetor no espaço $\vec{u} \in \R^3$
  \[
    \pause
    \vec{u}
    = ( u_1, u_2, u_3 )
    \pause
    = \vetortri{u_1}{u_2}{u_3}
  \]

  \pause
  Pelo Teorema de Pitágoras
  \[
    \pause
    \norm{\vec{u}}
    \pause
    = \sqrt{ u_1^2 + u_2^2 + u_3^2 }
  \]
\end{frame}

%------------------------------------------------------------------------------%
\begin{frame}[t]{Cálculo do Módulo}
  \vspace{1\baselineskip}

  De forma geral, $\vec{u} \in \R^n$
  \[
    \pause
    \vec{u}
    \pause
    = ( u_1, u_2, \dots, u_n )
    \pause
    = \begin{pmatrix}[c]
      u_1 \\
      u_2 \\
      \vdots \\
      u_n
    \end{pmatrix}
  \]
  \pause
  Pelo Teorema de Pitágoras
  \[
    \pause
    \norm{\vec{u}}
    \pause
    = \sqrt{ u_1^2 + u_2^2 + \cdots + u_n^2 }
    \pause
    = \sqrt{ \; \sum_{i=1}^{n} u_i^2 \; }
  \]
\end{frame}

%------------------------------------------------------------------------------%
\begin{frame}{Propriedades do Módulo}
  Para quaisquer vetores $\vec{u}$ e $\vec{v}$ de mesma dimensão
  \pause
  \[
    \norm{\vec{u}} \geq 0
  \]

  \pause
  \[
    \norm{\vec{u}} = 0
    \qquad\Leftrightarrow\qquad
    \vec{u} = \vec{0}
  \]

  \pause
  \[
    \norm{-\vec{u}} = \norm{\vec{u}}
  \]

  \pause
  \[
    \norm{\vec{u}+\vec{v}}
    \leq
    \norm{\vec{u}} + \norm{\vec{v}}
    \pause
    \qquad
    \text{\emph{desigualdade triangular}}
  \]
\end{frame}

%------------------------------------------------------------------------------%
%------------------------------------------------------------------------------%
\begin{question}
  Determine o módulo do vetor
  \(
    \vec{u} = ( 6, 8 )
  \)
\end{question}

%------------------------------------------------------------------------------%
\begin{resolution}{Solução}
  \begin{align*}
  \onslide<+->{\norm{\vec{u}}}
  \onslide<+->{& = \sqrt{ 6^2 + 8^2 } \\}
  \onslide<+->{& = \sqrt{ 36 + 64 } \\}
  \onslide<+->{& = \sqrt{ 100 } \\}
  \onslide<+->{& = 10}
  \end{align*}
\end{resolution}

%------------------------------------------------------------------------------%
\endinput

%------------------------------------------------------------------------------%
\begin{question}
  Determine o valor de $k$ para que o vetor
  \[
    \vec{u} = \vetor{k}{4}
  \]
  possua módulo igual a 5
\end{question}

%------------------------------------------------------------------------------%
\begin{resolution}{Solução}
  \begin{align*}
    \onslide<+->{\norm{\vec{u}} = \sqrt{ k^2 + 16 } & = 5                          \\}
    \onslide<+->{k^2 + 16                           & = 25                         \\}
    \onslide<+->{k^2                                & = 25 - 16                    \\}
    \onslide<+->{k^2                                & = 9                          \\}
    \onslide<+->{\abs{k}                            & =3                           \\}
    \onslide<+->{k                                  & =\pm3                        \\}
    \onslide<+->{k=3                                & \quad \text{ou}\quad k=-3}
  \end{align*}
\end{resolution}

%------------------------------------------------------------------------------%
\endinput


%------------------------------------------------------------------------------%
\section{Produto por Escalar}

%------------------------------------------------------------------------------%
\begin{frame}{Produto de um Vetor por um Escalar}
  O produto de $\vec{u} \in \R^n$
  \pause
  por $\lambda \in \R$ é um novo vetor
  \[
    \pause
    \vec{v} = \lambda \vec{u}
  \]

  \pause
  O vetor $\vec{u}$ é escalado pelo valor $\lambda$

  \pause
  O vetor $\vec{v}$ é colinear, ou paralelo, ao vetor $\vec{u}$
\end{frame}

%------------------------------------------------------------------------------%
\begin{frame}{Produto de um Vetor por um Escalar}
  \pause
  Se \; $\lambda>1$   \tabto{30mm} o vetor aumenta de comprimento

  \pause
  \smallskip
  Se \; $0<\lambda<1$ \tabto{30mm} o vetor diminui de comprimento

  \pause
  \smallskip
  Se \; $\lambda<0$   \tabto{30mm} o sentido é invertido

  \pause
  \smallskip
  Se \; $\lambda=0$   \tabto{30mm} obtemos o vetor nulo
\end{frame}

%------------------------------------------------------------------------------%
\framedgraphic*{Interpretação Geométrica}{E-02-Operacoes-prod-escalar-1}{}
\framedgraphic*{Interpretação Geométrica}{E-02-Operacoes-prod-escalar-2}{}
\framedgraphic*{Interpretação Geométrica}{E-02-Operacoes-prod-escalar-3}{}

%------------------------------------------------------------------------------%
\begin{frame}{Cálculo do Produto por Escalar}
  Para
  \(
    \lambda \in \R
  \)
  \pause
  e
  \(
    \vec{u}
    = ( u_1, u_2 )
    \in \R^2
  \)

  \pause
  \[
    \lambda\vec{u}
    \pause
    = \lambda \left( u_1, u_2 \right)
    \pause
    = \left( \lambda u_1, \lambda u_2 \right)
    \pause
    = \vetor{\lambda u_1}{\lambda u_2}
    \pause
    = \lambda \vetor{u_1}{u_2}
  \]
\end{frame}

%------------------------------------------------------------------------------%
%------------------------------------------------------------------------------%
\begin{question}
  Dado
  \[
    \vec{u} = ( 2, -3 )
  \]
  Determine
  \[
    3\vec{u}
    \qquad\text{e}\qquad
    -2\vec{u}
  \]
\end{question}

%------------------------------------------------------------------------------%
\begin{resolution}{Solução}
\begin{columns}
\column{0.3\linewidth}
  \begin{align*}
    \onslide<+->{\vec{v}}
    \onslide<+->{& = 3\vec{u}      \\}
    \onslide<+->{& = 3 ( 2, -3 )   \\}
    \onslide<+->{& = ( 6, -9 )          }
  \end{align*}
\column{0.5\linewidth}
  \begin{align*}
    \onslide<+->{\vec{w}}
    \onslide<+->{& = -2\vec{u}     \\}
    \onslide<+->{& = -2 ( 2, -3 )  \\}
    \onslide<+->{& = ( -4, 6 )}
  \end{align*}
\end{columns}
\end{resolution}

%------------------------------------------------------------------------------%
\endinput



%------------------------------------------------------------------------------%
\begin{frame}{Propriedades do Produto por Escalar}
  Para escalares $\lambda$ e $\mu$ e vetores $\vec{u}$ e $\vec{v}$
  \pause
  \[
    1 \, \vec{u} = \vec{u}
  \]

  \pause
  \[
    0 \, \vec{u} = \vec{0}
  \]

  \pause
  \[
    (-1)\vec{u} = -\vec{u}
  \]

  \pause
  \[
    \lambda(\mu\vec{u}) = (\lambda\mu)\vec{u}
  \]

  \pause
  \[
    (\lambda+\mu)\vec{u} = \lambda\vec{u}+\mu\vec{u}
  \]

  \pause
  \[
    \lambda(\vec{u}+\vec{v}) = \lambda\vec{u}+\lambda\vec{v}
  \]
\end{frame}

%------------------------------------------------------------------------------%
\begin{frame}{Relação Entre Módulo e Produto por Escalar}
  O módulo do produto por escalar satisfaz
  \[
    \norm{\lambda\vec{u}}
    =
    \abs{\lambda} \; \norm{\vec{u}}
  \]
\end{frame}

%------------------------------------------------------------------------------%
\input{ex/E-02-Operacoes-ex-04}
%------------------------------------------------------------------------------%
\begin{question}
  Determine $p$ para que os vetores \(\vec{u}\) e \(\vec{v}\) sejam paralelos
  \[
    \vec{u} = ( 2, -1, 3 )
  \]
  \[
    \vec{v} = (-6, 3, p)
  \]
\end{question}

%------------------------------------------------------------------------------%
\begin{resolution}{Solução}
  \onslide<+->{Os vetores serão paralelos se um for múltiplo do outro,}
  \onslide<+->{isto é, se \emph{existir $\lambda$} tal que}
  \begin{align*}
    \onslide<+->{\vec{v}    & = \lambda\vec{u}       \\}
    \onslide<+->{(-6, 3, p) & = \lambda ( 2, -1, 3 ) \\}
    \onslide<+->{(-6, 3, p) & = \big( \lambda 2,\, \lambda (-1),\, \lambda 3 \big)\\}
    \onslide<+->{(-6, 3, p) & = \big( 2\lambda,\, -\lambda,\, 3\lambda \big)}
  \end{align*}
\end{resolution}

%------------------------------------------------------------------------------%
\begin{resolution}{Solução}
  \begin{columns}[t]
  \column{0.25\linewidth}
    \begin{align*}
      \onslide<+->{-6       & = 2\lambda \\}
      \onslide<+->{2\lambda & = -6       \\}
      \onslide<+->{\lambda  & = -3}
    \end{align*}
  \column{0.4\linewidth}
    \begin{align*}
      \onslide<+->{3        & = -\lambda \\}
      \onslide<+->{-\lambda & = 3        \\}
      \onslide<+->{\lambda  & = -3}
    \end{align*}
    \qquad\onslide<+->{compatível}
  \column{0.35\linewidth}
    \begin{align*}
      \onslide<+->{p & = 3\lambda \\}
      \onslide<+->{  & = 3(-3)    \\}
      \onslide<+->{  & = -9}
    \end{align*}
  \end{columns}
\end{resolution}

%------------------------------------------------------------------------------%
\endinput


%------------------------------------------------------------------------------%
\section{Soma de vetores}

%------------------------------------------------------------------------------%
\begin{frame}{Soma de Vetores}
  Dados dois vetores
  \(
    \vec{u}
  \)
  e
  \(
    \vec{v}
  \)
  em \(\R^n\)

  \pause
  Representamos sua soma por
  \[
    \pause
    \vec{w} = \vec{u} + \vec{v}
  \]

  \pause
  A soma representa o ``deslocamento'' acumulado dos dois vetores
\end{frame}

%------------------------------------------------------------------------------%
\begin{frame}{Interpretação Geométrica}
  Para realizar a soma
  \(
    \vec{w} = \vec{u}+\vec{v}
  \)
  geometricamente
  \pause
  \begin{itemize}
      \item colocamos a origem de $\vec{v}$ na extremidade de $\vec{u}$
      \pause
      \item o vetor soma \(\vec{w}\) liga a origem de $\vec{u}$ à extremidade de $\vec{v}$
  \end{itemize}
\end{frame}

%------------------------------------------------------------------------------%
\framedgraphic*{Interpretação Geométrica}{E-02-Operacoes-soma-1}{}
\framedgraphic*{Interpretação Geométrica}{E-02-Operacoes-soma-2}{}
\framedgraphic*{Interpretação Geométrica}{E-02-Operacoes-soma-3}{}

%------------------------------------------------------------------------------%
\begin{frame}{Definição Algébrica da Soma}
  Dados
  \(
    \vec{u} = \vetor{u_1}{u_2}
  \)
  e
  \(
    \vec{v} = \vetor{v_1}{v_2}
  \)

  \pause
  \medskip
  A soma é definida componente a componente
  \pause
  \[
    \vec{u}+\vec{v}
    \pause
    \;=\; \vetor{u_1}{u_2} + \vetor{v_1}{v_2}
    \pause
    \;=\; \vetor{u_1+v_1}{u_2+v_2}
  \]
\end{frame}

%------------------------------------------------------------------------------%
%------------------------------------------------------------------------------%
\begin{question}
  Calcule a soma dos vetores
  \[
    \vec{u} = (2, 3)
  \]
  \[
    \vec{v} = (4, -1)
  \]

\end{question}

%------------------------------------------------------------------------------%
\begin{resolution}{Solução}
  \begin{align*}
    \onslide<+->{\vec{u}+\vec{v}}
    \onslide<+->{& = (2, 3) + (4, -1) \\}
    \onslide<+->{& = (2+4,\, 3-1) \\}
    \onslide<+->{& = (6, 2)}
  \end{align*}
\end{resolution}

%------------------------------------------------------------------------------%
\endinput


%------------------------------------------------------------------------------%
\begin{frame}{Propriedades da soma}
  Para quaisquer vetores $\vec{u}$, $\vec{v}$ e $\vec{w}$

  \pause
  \[
    \vec{u}+\vec{v} = \vec{v}+\vec{u}
  \]

  \pause
  \[
    (\vec{u}+\vec{v})+\vec{w} = \vec{u}+(\vec{v}+\vec{w})
  \]

  \pause
  \[
    \vec{u}+\vec{0} = \vec{u}
  \]

  \pause
  \[
    \vec{u}+(-\vec{u}) = \vec{0}
  \]
\end{frame}

%------------------------------------------------------------------------------%
\begin{frame}{Subtração de vetores}
  \onslide<+->{A subtração é definida pela soma com o vetor oposto}
  \begin{align*}
    \onslide<+->{\vec{u}-\vec{v}}
    \onslide<+->{& = \vec{u} + (-\vec{v})}
    \onslide<+->{  = \vec{u} + (-1)\vec{v}                   \\[3mm]}
    \onslide<+->{& = \vetor{u_1}{u_2} + (-1)\vetor{v_1}{v_2} \\[3mm]}
    \onslide<+->{& = \vetor{u_1}{u_2} + \vetor{-v_1}{-v_2}   \\[3mm]}
    \onslide<+->{& = \vetor{u_1-v_1}{u_2-v_2}}
  \end{align*}
\end{frame}

%%================================================
%\begin{frame}{Interpretação geométrica da diferença}
%
%A diferença
%
%\[
%\vec{u}-\vec{v}
%\]
%
%pode ser interpretada como o vetor que vai da extremidade de
% $\vec{v}$ à extremidade de $\vec{u}$ quando os dois
% vetores possuem a mesma origem
%\begin{center}
%\begin{asy}
%size(8cm,0);
%
%pair O=(0,0);
%pair A=(4,1);
%pair B=(1,3);
%
%draw(O--A,Arrow);
%draw(O--B,Arrow);
%draw(B--A,Arrow);
%
%dot(O);
%dot(A);
%dot(B);
%
%label("$O$",O,SW);
%label("$\vec{u}$",A/2,SE);
%label("$\vec{v}$",B/2,NW);
%label("$\vec{u}-\vec{v}$",(A+B)/2,NE);
%
%\end{asy}
%\end{center}
%\end{frame}

%------------------------------------------------------------------------------%
%------------------------------------------------------------------------------%
\begin{question}
  Dados
  \[
    \vec{u} = ( 3, -2 )
  \]
  \[
    \vec{v} = ( -1, 5 )
  \]
  Determine
  \[
    \vec{w} = \vec{u} + \vec{v}
  \]

\end{question}

%------------------------------------------------------------------------------%
\begin{resolution}{Solução}
  \begin{align*}
    \onslide<+->{\vec{w}}
    \onslide<+->{& = \vec{u} + \vec{v} \\}
    \onslide<+->{& = \vetor{3}{-2} + \vetor{-1}{5}  \\}
    \onslide<+->{& = \vetor{3+(-1)}{-2+5} \\}
    \onslide<+->{& = \vetor{2}{3}}
  \end{align*}
\end{resolution}

%------------------------------------------------------------------------------%
\endinput


%------------------------------------------------------------------------------%
\begin{question}
  Dados
  \[
    \vec{u} = (3, 4)
  \]
  \[
    \vec{v} = (4, -3)
  \]
  Determine o módulo da soma dos vetores
\end{question}

%------------------------------------------------------------------------------%
\begin{resolution}{Solução}
\begin{columns}[t]
\column{0.5\linewidth}
  \onslide<+->{Calculando a Soma}
  \begin{align*}
    \onslide<+->{\vec{u} + \vec{v}}
    \onslide<+->{& = (3, 4) + (4, -3) \\}
    \onslide<+->{& = (3+4,\; 4-3) \\}
    \onslide<+->{& = (7, 1)}
  \end{align*}
\column{0.5\linewidth}
  \onslide<+->{Calculando o módulo}
  \begin{align*}
    \onslide<+->{\norm{\vec{u}+\vec{v}}}
    \onslide<+->{& = \norm{(7, 1)} \\}
    \onslide<+->{& = \sqrt{7^2+1^2}\\}
    \onslide<+->{& = \sqrt{50}\\}
    \onslide<+->{& = 5\sqrt{2}}
  \end{align*}
\end{columns}
\end{resolution}

%------------------------------------------------------------------------------%
\endinput

%------------------------------------------------------------------------------%
\begin{question}
  Dados
  \[
    \vec{u} = (1, 2)
  \]
  \[
    \vec{v} = (3, -1)
  \]
  Determine
  \[
    \vec{w} = 2\vec{u} - 3\vec{v}
  \]
\end{question}

%------------------------------------------------------------------------------%
\begin{resolution}{Solução}
  \[
    2\vec{u}
    \pause
    \;=\; 2(1, 2)
    \pause
    \;=\; (2\times 1, \; 2\times 2)
    \pause
    \;=\; (2, 4)
  \]

  \pause
  \[
    3\vec{v}
    \pause
    \;=\; 3(3, -1)
    \pause
    \;=\; (3\time 3, \; 3(-1))
    \pause
    \;=\; (9, -3)
  \]

  \pause
  Portanto
  \[
    \pause
    \vec{w}
    \pause
    \;=\; 2\vec{u}-3\vec{v}
    \pause
    \;=\; (2, 4) - (9, -3)
    \pause
    \;=\; (-7, 7)
  \]
\end{resolution}

%------------------------------------------------------------------------------%
\endinput

%------------------------------------------------------------------------------%
\begin{question}
  Avalie
  \(
    \norm{\vec{u} + 2 \vec{v}}
  \)
  para
  \[
    \vec{u} = \vetortri{2}{-1}{3}
  \]
  \[
    \vec{v} = \vetortri{1}{4}{-2}
  \]
\end{question}

%------------------------------------------------------------------------------%
\begin{resolution}{Solução}
  \[
    2\vec{v}
    \pause
    \;=\; 2 \vetortri{2}{-1}{3}
    \pause
    \;=\; \vetortri{2\times 2}{2(-1)}{2\times 3}
    \pause
    \;=\; \vetortri{2}{8}{-4}
  \]
  \pause
  \[
    \vec{u} + 2 \vec{v}
    \pause
    \;=\; \vetortri{2}{-1}{3} + \vetortri{2}{8}{-4}
    \pause
    \;=\; \vetortri{4}{7}{-1}
  \]
  \pause
  \[
    \norm{\vec{u} + 2 \vec{v}}
    \pause
    \;=\; \norm{\vetortri{4}{7}{-1}}
    \pause
    \;=\; \sqrt{4^2 + 7^2 + (-1)^2\,}
    \pause
    \;=\; \sqrt{16 + 49 + 1\,}
    \pause
    \;=\; \sqrt{66}
  \]
\end{resolution}

%------------------------------------------------------------------------------%
\endinput



%------------------------------------------------------------------------------%
\begin{question}
  Dados
  \(
    \vec{u} = \vetor{2}{-1}
  \) e
  \(
    \vec{v} = \vetor{-4}{3}
  \)

  \bigskip
  Determine $\lambda$ de modo que
  \(
    \norm{2\vec{u} + \vec{v} + \vetor{\lambda}{1}} = 3
  \)
\end{question}

%------------------------------------------------------------------------------%
\begin{resolution}{Solução}
  \begin{align*}
    \onslide<+->{\vec{w}}
    \onslide<+->{& = 2\vec{u} + \vec{v} + \vetor{\lambda}{1}              \\}
    \onslide<+->{& = 2\vetor{2}{-1} + \vetor{-4}{3} + \vetor{\lambda}{1}  \\}
    \onslide<+->{& = \vetor{4}{-2} + \vetor{-4}{3} + \vetor{\lambda}{1}   \\}
    \onslide<+->{& = \vetor{4-4+\lambda}{-2+3+1}                          \\}
    \onslide<+->{& = \vetor{\lambda}{2}}
  \end{align*}
\end{resolution}

%------------------------------------------------------------------------------%
\begin{resolution}{Solução}
  \begin{align*}
    \onslide<+->{\norm{\vec{w}}}
    \onslide<+->{= \norm{\vetor{\lambda}{2}}}
    \onslide<+->{& = 3 \\}
    \onslide<+->{\sqrt{\lambda^2 + 2^2} & = 3        \\}
    \onslide<+->{\lambda^2 + 4          & = 3^2      \\}
    \onslide<+->{\lambda^2              & = 9 - 4    \\}
    \onslide<+->{\lambda^2              & = 5        \\}
    \onslide<+->{\abs{\lambda}          & = \sqrt{5} \\}
    \onslide<+->{\lambda = \sqrt{5}     & }
    \onslide<+->{\quad \text{ou} \quad \lambda = -\sqrt{5}}
  \end{align*}
\end{resolution}

%------------------------------------------------------------------------------%
\endinput


%------------------------------------------------------------------------------%
%\section{Lista Mínima}
%
%%------------------------------------------------------------------------------%
%\begin{frame}{Lista Mínima}
%  \begin{enumerate}
%    \item
%  \end{enumerate}
%
%  \vfill
%  Atenção: A prova é baseada no livro, não nas apresentações
%\end{frame}

%------------------------------------------------------------------------------%
\end{document}
%------------------------------------------------------------------------------%


